\section*{Changes since Assignment 3}
The feedback accepted the matched candidates as an exploratory improvement and asked us to confirm the
main conclusions under a protocol fixed in advance. We wrote that protocol, tagged it in our repository before any run at a new seed, and reran the stability comparison at ten seeds on the nine small datasets and five
on PBMC14K. Every run now draws 100 subsamples, and the first 25 reproduce the original protocol exactly,
so one run gives both. We added PCA plus $k$-means as a textbook baseline under the same rule, and
scored every run with chooseR's own R code as well as ours, which showed a third difference besides the
number of subsamples and the bootstrap type: chooseR resamples clusters where we resampled cells. We also explain why matching changes the ranking (Section~\ref{sec:why}), test whether the setting that stands for a cluster count matters (Section~\ref{sec:variants}), and show how each lead is spread across the datasets (Section~\ref{sec:confirm}). Section~\ref{sec:design} says which secondary results keep three seeds or one, and why. The paper is reframed around
how label-free selection changes the ranking of clustering methods and when such a comparison is
reliable, with MICA against Leiden as the case that exposed it. The graph comparison moves to the
appendix, and the large-scale dataset the course requires is planned for the November 30 draft.
